% 90-120 minute graduate lecture: adapting pretrained classifiers, fine-grained recognition,
% deep metric learning for retrieval and re-identification, and open-set recognition.
% Narrative thread: a closed-set classifier first gets sharper (adaptation, fine-grained), then
% learns a similarity instead of a label set (metric learning, retrieval, re-ID), and finally
% learns to say "none of the above" (open-set, OOD detection, category discovery).
\documentclass[10pt, aspectratio=169]{beamer}

\usepackage{tikz}
\usetikzlibrary{positioning, arrows.meta, shapes.geometric, calc, fit, backgrounds}

\input{preamble/packages}
\input{preamble/commands}
\input{preamble/beamer_settings}

\title[Image Classification \& Recognition]{Image Classification and Recognition}
\subtitle{Fine-Grained Recognition, Metric Learning, and Open-Set Recognition}

\begin{document}

\input{sections/cover}
\input{sections/image-recognition/toc}
\input{sections/image-recognition/learning_outcomes}
\input{sections/image-recognition/introduction}
\input{sections/image-recognition/classification_adaptation}
\input{sections/image-recognition/fine_grained}
\input{sections/image-recognition/metric_learning}
\input{sections/image-recognition/retrieval}
\input{sections/image-recognition/reidentification}
\input{sections/image-recognition/open_set}
\input{sections/image-recognition/open_world}
\input{sections/image-recognition/limitations}
\input{sections/image-recognition/summary}
\input{sections/image-recognition/references}

\end{document}
